% GENERATED FILE. Edit canonical lessons, then run npm run generate:books.
% canonical-source-hash: fnv1a64:ef8aef8aaaf1b030
% canonical-lessons: PA-C276-mull, PA-C276-bhaa, PA-C276-ghatauna, PA-C276-sauda, PA-C276-mufat

\chapter{What Is the Price?}
\label{ch:pa-a2-what-is-the-price}
\hlchaptermodality{276}

\begin{chapteropening}
\textbf{By the end of this chapter:} \emph{I can ask the price and ask for less.}

Complete the last lesson of chapter 276: I can ask the price and ask for less.
\end{chapteropening}

\section[\texorpdfstring{\pa{ਮੁੱਲ} (mull)}{ਮੁੱਲ (mull)}]{\pa{ਮੁੱਲ} (mull) — a price}

\label{lesson:PA-C276-mull}



\begin{quote}
Before the new one: say the Punjabi for a dozen, then the Punjabi for a bag.
\end{quote}

\subsection*{You'll want to know: \pa{ਮੁੱਲ}}
\textbf{\pa{ਮੁੱਲ}} — \emph{mull} — \textquotedblleft{}a price\textquotedblright{}.

\begin{quote}
\pa{ਕਮੀਜ਼} \pa{ਦਾ} \pa{ਮੁੱਲ} \pa{ਕੀ} \pa{ਹੈ}?
\end{quote}

\begin{quote}
\emph{kamīz dā mull kī hai?}
\end{quote}

\begin{quote}
What is the price of the shirt?
\end{quote}

A masculine noun.

\begin{cousinweb}
\textbf{\pa{ਮੁੱਲ}} (\emph{mull}) is inherited from Sanskrit \emph{mūlya}, \textquotedblleft{}original value, price\textquotedblright{}; CDIAL lists Punjabi \emph{mull}, \textquotedblleft{}price\textquotedblright{}, under it (CDIAL 10257).
\end{cousinweb}

\begin{grammarlens}[title={What you've built}]
One more word for this chapter.
\end{grammarlens}

\subsection*{Your turn}
\begin{itemize}
\raggedright
  \item \emph{Say it:} \emph{mull}
  \item \emph{Say it:} \emph{mull}, once more
  \item \emph{Say it:} \emph{kamīz dā mull kī hai?} — What is the price of the shirt?
  \item \emph{Recall:} say the Punjabi for a dozen, then the Punjabi for a bag, then say \emph{mull} again
\end{itemize}

\begin{tcolorbox}[breakable,colback=teal!4,colframe=teal!35!black,title={Before you move on}]
What does \pa{ਮੁੱਲ} mean? (\textquotedblleft{}A price\textquotedblright{}.) Say it once more.
\end{tcolorbox}
\section[\texorpdfstring{\pa{ਭਾਅ} (bhāa)}{ਭਾਅ (bhāa)}]{\pa{ਭਾਅ} (bhāa) — a rate; the market rate}

\label{lesson:PA-C276-bhaa}



\begin{quote}
Before the new one: say the Punjabi for a bag, then the Punjabi for a price.
\end{quote}

\subsection*{You'll want to know: \pa{ਭਾਅ}}
\textbf{\pa{ਭਾਅ}} — \emph{bhāa} — \textquotedblleft{}a rate; the market rate\textquotedblright{}.

\begin{quote}
\pa{ਅੱਜ} \pa{ਸੇਬਾਂ} \pa{ਦਾ} \pa{ਭਾਅ} \pa{ਕੀ} \pa{ਹੈ}?
\end{quote}

\begin{quote}
\emph{ajj sebāṁ dā bhāa kī hai?}
\end{quote}

\begin{quote}
What is the rate for apples today?
\end{quote}

A masculine noun.

\begin{cousinweb}
\textbf{\pa{ਭਾਅ}} (\emph{bhāa}) is filed by CDIAL under Sanskrit \emph{bhāva}, \textquotedblleft{}being, state\textquotedblright{}; CDIAL lists Punjabi \emph{bhāu}, \textquotedblleft{}market rate\textquotedblright{}, there (CDIAL 9475).
\end{cousinweb}

\begin{grammarlens}[title={What you've built}]
One more word for this chapter.
\end{grammarlens}

\subsection*{Your turn}
\begin{itemize}
\raggedright
  \item \emph{Say it:} \emph{bhāa}
  \item \emph{Say it:} \emph{bhāa}, once more
  \item \emph{Say it:} \emph{ajj sebāṁ dā bhāa kī hai?} — What is the rate for apples today?
  \item \emph{Recall:} say the Punjabi for a bag, then the Punjabi for a price, then say \emph{bhāa} again
  \item \emph{Write it:} copy \textbf{\pa{ਭਾਅ}} once, checking the whole word
\end{itemize}

\begin{tcolorbox}[breakable,colback=teal!4,colframe=teal!35!black,title={Before you move on}]
What does \pa{ਭਾਅ} mean? (\textquotedblleft{}A rate; the market rate\textquotedblright{}.) Say it once more.
\end{tcolorbox}
\section[\texorpdfstring{\pa{ਘਟਾਉਣਾ} (ghaṭāuṇā)}{ਘਟਾਉਣਾ (ghaṭāuṇā)}]{\pa{ਘਟਾਉਣਾ} (ghaṭāuṇā) — to lower; to reduce}

\label{lesson:PA-C276-ghatauna}



\begin{quote}
Before the new one: say the Punjabi for a price, then the Punjabi for a rate.
\end{quote}

\subsection*{You'll want to know: \pa{ਘਟਾਉਣਾ}}
\textbf{\pa{ਘਟਾਉਣਾ}} — \emph{ghaṭāuṇā} — \textquotedblleft{}to lower; to reduce\textquotedblright{}.

\begin{quote}
\pa{ਇਹ} \pa{ਦੁਕਾਨ} \pa{ਮੁੱਲ} \pa{ਨਹੀਂ} \pa{ਘਟਾਉਂਦੀ} \pa{ਹੈ।}
\end{quote}

\begin{quote}
\emph{eh dukān mull nahīṁ ghaṭāundī hai.}
\end{quote}

\begin{quote}
This shop does not lower its prices.
\end{quote}

Past (I …): \pa{ਮੈਂ} \pa{ਮੁੱਲ} \pa{ਘਟਾਇਆ}, \emph{maiṁ mull ghaṭāiā}.

\begin{cousinweb}
\textbf{\pa{ਘਟਾਉਣਾ}} (\emph{ghaṭāuṇā}) is the causative partner of \emph{ghaṭṇā}, \textquotedblleft{}to decrease\textquotedblright{}; CDIAL lists Punjabi \emph{ghaṭṇā} under a reconstructed \emph{*ghaṭṭati}, \textquotedblleft{}decreases, is wanting\textquotedblright{} (CDIAL 4415).
\end{cousinweb}

\begin{grammarlens}[title={What you've built}]
One more word for this chapter.
\end{grammarlens}

\subsection*{Your turn}
\begin{itemize}
\raggedright
  \item \emph{Say it:} \emph{ghaṭāuṇā}
  \item \emph{Say it:} \emph{ghaṭāuṇā}, once more
  \item \emph{Say it:} \emph{eh dukān mull nahīṁ ghaṭāundī hai.} — This shop does not lower its prices.
  \item \emph{Recall:} say the Punjabi for a price, then the Punjabi for a rate, then say \emph{ghaṭāuṇā} again
  \item \emph{Write it:} copy \textbf{\pa{ਘਟਾਉਣਾ}} once, checking the whole word
\end{itemize}

\begin{tcolorbox}[breakable,colback=teal!4,colframe=teal!35!black,title={Before you move on}]
What does \pa{ਘਟਾਉਣਾ} mean? (\textquotedblleft{}To lower; to reduce\textquotedblright{}.) Say it once more.
\end{tcolorbox}
\section[\texorpdfstring{\pa{ਸੌਦਾ} (saudā)}{ਸੌਦਾ (saudā)}]{\pa{ਸੌਦਾ} (saudā) — a purchase; a deal}

\label{lesson:PA-C276-sauda}



\begin{quote}
Before the new one: say the Punjabi for a rate, then the Punjabi for to lower.
\end{quote}

\subsection*{You'll want to know: \pa{ਸੌਦਾ}}
\textbf{\pa{ਸੌਦਾ}} — \emph{saudā} — \textquotedblleft{}a purchase; a deal\textquotedblright{}.

\begin{quote}
\pa{ਇਹ} \pa{ਚੰਗਾ} \pa{ਸੌਦਾ} \pa{ਹੈ।}
\end{quote}

\begin{quote}
\emph{eh changā saudā hai.}
\end{quote}

\begin{quote}
This is a good deal.
\end{quote}

A masculine noun.

\begin{cousinweb}
\textbf{\pa{ਸੌਦਾ}} (\emph{saudā}) is a loan, usually traced to Persian \emph{saudā}, \textquotedblleft{}trade, a bargain\textquotedblright{}.
\end{cousinweb}

\begin{grammarlens}[title={What you've built}]
One more word for this chapter.
\end{grammarlens}

\subsection*{Your turn}
\begin{itemize}
\raggedright
  \item \emph{Say it:} \emph{saudā}
  \item \emph{Say it:} \emph{saudā}, once more
  \item \emph{Say it:} \emph{eh changā saudā hai.} — This is a good deal.
  \item \emph{Recall:} say the Punjabi for a rate, then the Punjabi for to lower, then say \emph{saudā} again
  \item \emph{Write it:} copy \textbf{\pa{ਸੌਦਾ}} once, checking the whole word
\end{itemize}

\begin{tcolorbox}[breakable,colback=teal!4,colframe=teal!35!black,title={Before you move on}]
What does \pa{ਸੌਦਾ} mean? (\textquotedblleft{}A purchase; a deal\textquotedblright{}.) Say it once more.
\end{tcolorbox}
\section[\texorpdfstring{\pa{ਮੁਫ਼ਤ} (mufat)}{ਮੁਫ਼ਤ (mufat)}]{\pa{ਮੁਫ਼ਤ} (mufat) — free of charge}

\label{lesson:PA-C276-mufat}



\begin{quote}
Before the new one: say the Punjabi for to lower, then the Punjabi for a purchase.
\end{quote}

\subsection*{You'll want to know: \pa{ਮੁਫ਼ਤ}}
\textbf{\pa{ਮੁਫ਼ਤ}} — \emph{mufat} — \textquotedblleft{}free of charge\textquotedblright{}.

\begin{quote}
\pa{ਦੁਕਾਨ} \pa{ਵਿੱਚ} \pa{ਥੈਲਾ} \pa{ਮੁਫ਼ਤ} \pa{ਹੈ।}
\end{quote}

\begin{quote}
\emph{dukān vich thailā mufat hai.}
\end{quote}

\begin{quote}
The bag is free in the shop.
\end{quote}

\begin{cousinweb}
\textbf{\pa{ਮੁਫ਼ਤ}} (\emph{mufat}) is a loan, usually traced to Persian \emph{muft}, \textquotedblleft{}free, for nothing\textquotedblright{}.
\end{cousinweb}

\begin{grammarlens}[title={What you've built}]
One more word for this chapter.
\end{grammarlens}

\subsection*{Your turn}
\begin{itemize}
\raggedright
  \item \emph{Say it:} \emph{mufat}
  \item \emph{Say it:} \emph{mufat}, once more
  \item \emph{Say it:} \emph{dukān vich thailā mufat hai.} — The bag is free in the shop.
  \item \emph{Recall:} say the Punjabi for to lower, then the Punjabi for a purchase, then say \emph{mufat} again
  \item \emph{Write it:} copy \textbf{\pa{ਮੁਫ਼ਤ}} once, checking the whole word
\end{itemize}

\begin{tcolorbox}[breakable,colback=teal!4,colframe=teal!35!black,title={Before you move on}]
What does \pa{ਮੁਫ਼ਤ} mean? (\textquotedblleft{}Free of charge\textquotedblright{}.) Say it once more.
\end{tcolorbox}
